\documentclass[11pt]{article}
\usepackage{amsmath, amssymb}
\usepackage{geometry}
\usepackage{array}
\usepackage{booktabs}
\geometry{margin=0.7in}
\setlength{\parindent}{0pt}
\setlength{\parskip}{0.35em}
\renewcommand{\arraystretch}{1.15}

\newcommand{\cheatsheettitle}[2]{%
\begin{center}
{\LARGE #1}\\[0.4em]
{\large #2}
\end{center}
}


\begin{document}
\cheatsheettitle{Algebra II Complex Numbers and Polar Cheatsheet}{Module 12}

\section*{Complex Basics}
\[
i^2=-1,\qquad i^3=-i,\qquad i^4=1
\]
\[
\sqrt{-a}=i\sqrt a
\]
Standard form:
\[
a+bi
\]
Conjugate:
\[
a-bi
\]
\[
(a+bi)(a-bi)=a^2+b^2
\]

\section*{Complex Arithmetic}
\[
(a+bi)+(c+di)=(a+c)+(b+d)i
\]
\[
(a+bi)(c+di)=(ac-bd)+(ad+bc)i
\]
Divide by multiplying by the conjugate of the denominator.

\section*{Polar Coordinates}
\[
x=r\cos\theta,\qquad y=r\sin\theta
\]
\[
r^2=x^2+y^2,\qquad \tan\theta=\frac yx
\]
Choose $\theta$ using the quadrant of $(x,y)$.

\section*{Complex Polar Form}
\[
z=a+bi
\]
\[
r=\sqrt{a^2+b^2},\qquad \tan\theta=\frac ba
\]
\[
z=r(\cos\theta+i\sin\theta)=r\,\operatorname{cis}\theta
\]

\section*{Multiplication and Division in Polar Form}
\[
r_1\operatorname{cis}\theta_1\cdot r_2\operatorname{cis}\theta_2
=r_1r_2\operatorname{cis}(\theta_1+\theta_2)
\]
\[
\frac{r_1\operatorname{cis}\theta_1}{r_2\operatorname{cis}\theta_2}
=\frac{r_1}{r_2}\operatorname{cis}(\theta_1-\theta_2)
\]

\section*{De Moivre}
\[
\left(r\operatorname{cis}\theta\right)^n=r^n\operatorname{cis}(n\theta)
\]

\section*{Nth Roots}
For
\[
z=r\operatorname{cis}\theta
\]
the $n$ roots are:
\[
z_k=r^{1/n}\operatorname{cis}\left(\frac{\theta+2\pi k}{n}\right),
\qquad k=0,1,\dots,n-1
\]

\section*{Quick Checks}
\begin{itemize}
  \item Complex conjugates are useful for division and real-coefficient polynomial roots.
  \item Polar angles are not unique.
  \item Negative $r$ points in the opposite direction.
  \item Roots are evenly spaced around a circle.
\end{itemize}

\end{document}
